\section*{September 7, 2026: Transaction weights and the choice of price scale}
\textit{Codex draft; weighting decision confirmed by Jacob, proposed diagnostics below.}

Jacob chose equal weight per transaction for the working quarter-mile
comparison of 30 closed and 49 still-open schools. Giving each school area
equal weight within a year can give a few sales in a sparse area as much
influence as many sales elsewhere. Equal-site results remain a sensitivity
check. The existing transaction-weighted estimates already implement this
choice; the sample and model formulas do not change. Standard errors remain
clustered by school site, since sales near the same school are not independent.

The difference between dollar and log estimates needs an economic explanation.
A dollar model assumes that, without closure, the groups would have experienced
parallel changes in average dollar prices, conditional on the included
controls. A log model makes that assumption for average log prices.
Roth and Sant'Anna (2023) show that validity across outcome transformations
requires stronger restrictions than parallel trends for one chosen outcome.
Roth (2022) also cautions against selecting a specification because its
pre-treatment test fails to reject.

An illustrative example makes the distinction concrete. Suppose treated
homes initially cost \$200,000 and control homes \$100,000, and every home
appreciates by 10 percent without any treatment effect. The dollar difference
in differences is then positive \$10,000, while the log difference in
differences is zero. This is an illustration, not an explanation established
for these Chicago data. Differences in the upper tail or in which homes sell
could also separate average dollar prices from average log prices.

Codex recommends keeping transaction weights fixed while examining price
percentiles, expensive sales' contribution to annual means, and estimates
omitting one school site at a time. Property-type composition should be
inspected alongside prices. These checks would help distinguish a broadly
shared price change from concentration in a few sales or neighborhoods.
They do not authorize new price exclusions or choose logs over dollars.
A log-link model for mean dollar prices could provide a further comparison,
but it would require its own assumptions about the untreated price trend.

{\small Sources reviewed September 7, 2026: Roth and Sant'Anna,
``When Is Parallel Trends Sensitive to Functional Form?,''
\textit{Econometrica} 91 (2023), 737--747,
DOI: 10.3982/ECTA19402; Roth, ``Pretest with Caution,''
\textit{AER: Insights} 4 (2022), 305--322, DOI: 10.1257/aeri.20210236;
Solon, Haider, and Wooldridge, ``What Are We Weighting For?,''
\textit{Journal of Human Resources} 50 (2015), 301--316.
The weighting paper distinguishes population targets, precision, and sampling
as different reasons to weight observations.}

{\small Reproduction: the existing transaction-weighted figures are
\texttt{event\_studies.png} and \texttt{event\_studies\_semiannual.png} in
the \texttt{home\_prices\_twfe} audit's output directory. Its Makefile reproduces
the estimates. This entry records a weighting choice and literature review;
it introduces no new estimates or data vintage.}
